\section{Maximum Flow --- Edmonds-Karp Algorithm}

\subsection{Problem}

Given a directed graph $G = (V, E)$ with source $s$ and sink $t$, each edge $(u,v)$ has capacity $c(u,v)$. Find the maximum flow from $s$ to $t$ respecting capacity constraints.

$$\text{maximize} \quad \sum_{v} f(s,v) \quad \text{subject to} \quad 0 \le f(u,v) \le c(u,v)$$

\textbf{Edmonds-Karp} = Ford-Fulkerson with BFS for shortest augmenting paths. Guarantees $O(VE^2)$.

\subsection{Algorithm Walkthrough}

\begin{animation}[id="maxflow-ek", label="Edmonds-Karp Max Flow"]
\shape{G}{Graph}{
  nodes=["S","A","B","C","D","T"],
  edges=[("S","A"),("S","B"),("A","B"),("A","C"),("B","D"),("C","D"),("C","T"),("D","T")],
  directed=true,
  layout="hierarchical"
}
\shape{vars}{VariableWatch}{names=["total_flow","path","bottleneck"], label="State"}

% Initialize all edges with 0/cap labels
\apply{G.edge[(S,A)]}{value="0/10"}
\apply{G.edge[(S,B)]}{value="0/8"}
\apply{G.edge[(A,B)]}{value="0/5"}
\apply{G.edge[(A,C)]}{value="0/7"}
\apply{G.edge[(B,D)]}{value="0/10"}
\apply{G.edge[(C,D)]}{value="0/3"}
\apply{G.edge[(C,T)]}{value="0/9"}
\apply{G.edge[(D,T)]}{value="0/12"}

\step
\recolor{G.node[S]}{state=current}
\recolor{G.node[T]}{state=good}
\apply{vars.var[total_flow]}{value="0"}
\apply{vars.var[path]}{value="---"}
\apply{vars.var[bottleneck]}{value="---"}
\narrate{Flow network: each edge shows \textbf{flow/capacity}. All flows start at 0. Goal: maximize flow from $S$ to $T$.}

\step
\narrate{\textbf{Iteration 1:} BFS from $S$ finds shortest path $S \to A \to C \to T$.}
\recolor{G.edge[(S,A)]}{state=current}
\recolor{G.edge[(A,C)]}{state=current}
\recolor{G.edge[(C,T)]}{state=current}
\recolor{G.node[A]}{state=current}
\recolor{G.node[C]}{state=current}
\apply{vars.var[path]}{value="S-A-C-T"}

\step
\narrate{Bottleneck $= \min(10, 7, 9) = 7$. Edge $A \to C$ (cap 7) is the bottleneck.}
\apply{vars.var[bottleneck]}{value="7"}
\recolor{G.edge[(S,A)]}{state=good}
\recolor{G.edge[(A,C)]}{state=good}
\recolor{G.edge[(C,T)]}{state=good}

\step
\narrate{Push 7 units. Update flow on each edge. $A \to C$ is now saturated (7/7).}
\apply{G.edge[(S,A)]}{value="7/10"}
\apply{G.edge[(A,C)]}{value="7/7"}
\apply{G.edge[(C,T)]}{value="7/9"}
\recolor{G.edge[(A,C)]}{state=error}
\recolor{G.edge[(S,A)]}{state=done}
\recolor{G.edge[(C,T)]}{state=done}
\recolor{G.node[A]}{state=idle}
\recolor{G.node[C]}{state=idle}
\apply{vars.var[total_flow]}{value="7"}

\step
\narrate{\textbf{Iteration 2:} BFS. $A \to C$ saturated, so BFS finds $S \to B \to D \to T$.}
\apply{vars.var[path]}{value="S-B-D-T"}
\apply{vars.var[bottleneck]}{value="---"}
\recolor{G.edge[(S,B)]}{state=current}
\recolor{G.edge[(B,D)]}{state=current}
\recolor{G.edge[(D,T)]}{state=current}
\recolor{G.node[B]}{state=current}
\recolor{G.node[D]}{state=current}
\recolor{G.edge[(A,C)]}{state=dim}

\step
\narrate{Bottleneck $= \min(8, 10, 12) = 8$. Edge $S \to B$ (cap 8) is the bottleneck.}
\apply{vars.var[bottleneck]}{value="8"}
\recolor{G.edge[(S,B)]}{state=good}
\recolor{G.edge[(B,D)]}{state=good}
\recolor{G.edge[(D,T)]}{state=good}

\step
\narrate{Push 8 units. $S \to B$ saturated (8/8).}
\apply{G.edge[(S,B)]}{value="8/8"}
\apply{G.edge[(B,D)]}{value="8/10"}
\apply{G.edge[(D,T)]}{value="8/12"}
\recolor{G.edge[(S,B)]}{state=error}
\recolor{G.edge[(B,D)]}{state=done}
\recolor{G.edge[(D,T)]}{state=done}
\recolor{G.node[B]}{state=idle}
\recolor{G.node[D]}{state=idle}
\apply{vars.var[total_flow]}{value="15"}

\step
\narrate{\textbf{Iteration 3:} $S \to B$ saturated. Only $S \to A$ (residual $10-7=3$) leaves $S$. BFS: $S \to A \to B \to D \to T$.}
\apply{vars.var[path]}{value="S-A-B-D-T"}
\apply{vars.var[bottleneck]}{value="---"}
\recolor{G.edge[(S,A)]}{state=current}
\recolor{G.edge[(A,B)]}{state=current}
\recolor{G.edge[(B,D)]}{state=current}
\recolor{G.edge[(D,T)]}{state=current}
\recolor{G.node[A]}{state=current}
\recolor{G.node[B]}{state=current}
\recolor{G.node[D]}{state=current}
\recolor{G.edge[(S,B)]}{state=dim}

\step
\narrate{Residual: $S \to A$: 3, $A \to B$: 5, $B \to D$: 2, $D \to T$: 4. Bottleneck $= 2$ at $B \to D$.}
\apply{vars.var[bottleneck]}{value="2"}
\recolor{G.edge[(S,A)]}{state=good}
\recolor{G.edge[(A,B)]}{state=good}
\recolor{G.edge[(B,D)]}{state=good}
\recolor{G.edge[(D,T)]}{state=good}

\step
\narrate{Push 2 units. $B \to D$ now saturated (10/10).}
\apply{G.edge[(S,A)]}{value="9/10"}
\apply{G.edge[(A,B)]}{value="2/5"}
\apply{G.edge[(B,D)]}{value="10/10"}
\apply{G.edge[(D,T)]}{value="10/12"}
\recolor{G.edge[(B,D)]}{state=error}
\recolor{G.edge[(S,A)]}{state=done}
\recolor{G.edge[(A,B)]}{state=done}
\recolor{G.edge[(D,T)]}{state=done}
\recolor{G.node[A]}{state=idle}
\recolor{G.node[B]}{state=idle}
\recolor{G.node[D]}{state=idle}
\apply{vars.var[total_flow]}{value="17"}

\step
\narrate{\textbf{Iteration 4:} $S \to A$ has 1 residual. $A \to C$ saturated, $A \to B$ has 3. But $B \to D$ saturated --- dead end via $B$. Need \textbf{residual graph}: reverse edge $C \to A$ has 7 capacity (undo $A \to C$ flow).}
\apply{vars.var[path]}{value="searching..."}
\apply{vars.var[bottleneck]}{value="---"}
\recolor{G.edge[(S,A)]}{state=current}
\recolor{G.node[A]}{state=current}
\recolor{G.edge[(A,C)]}{state=highlight}
\recolor{G.edge[(B,D)]}{state=dim}

\step
\narrate{Residual path found: $S \to A$ (1) $\to$ \textbf{reverse} $A \to C$ (7) $\to C \to D$ (3) $\to D \to T$ (2). Bottleneck $= \min(1,7,3,2) = 1$.}
\apply{vars.var[path]}{value="S-A-[rev]C-D-T"}
\apply{vars.var[bottleneck]}{value="1"}
\recolor{G.edge[(S,A)]}{state=good}
\recolor{G.edge[(A,C)]}{state=path}
\recolor{G.edge[(C,D)]}{state=good}
\recolor{G.edge[(D,T)]}{state=good}
\recolor{G.node[C]}{state=good}
\recolor{G.node[D]}{state=good}
\annotate{G.edge[(A,C)]}{label="reverse!", color=warn}

\step
\narrate{Push 1 unit via residual. $A \to C$ flow \textbf{decreases} $7 \to 6$ (reverse edge used). $S \to A$: 10/10 (saturated!). Total $= 18$.}
\apply{G.edge[(S,A)]}{value="10/10"}
\apply{G.edge[(A,C)]}{value="6/7"}
\apply{G.edge[(C,D)]}{value="1/3"}
\apply{G.edge[(D,T)]}{value="11/12"}
\recolor{G.edge[(S,A)]}{state=error}
\recolor{G.edge[(A,C)]}{state=done}
\recolor{G.edge[(C,D)]}{state=done}
\recolor{G.edge[(D,T)]}{state=done}
\recolor{G.node[A]}{state=idle}
\recolor{G.node[C]}{state=idle}
\recolor{G.node[D]}{state=idle}
\apply{vars.var[total_flow]}{value="18"}

\step
\narrate{\textbf{No more paths:} $S \to A$ (10/10) and $S \to B$ (8/8) both saturated. BFS from $S$ finds nothing. Algorithm terminates.}
\apply{vars.var[path]}{value="none"}
\apply{vars.var[bottleneck]}{value="---"}
\recolor{G.edge[(S,A)]}{state=dim}
\recolor{G.edge[(S,B)]}{state=dim}

\step
\narrate{\textbf{Result:} Max flow $= 18$. Min-cut: $\{S\}$ vs $\{A,B,C,D,T\}$, capacity $10+8 = 18$.}
\recolor{G.node[S]}{state=error}
\recolor{G.node[A]}{state=done}
\recolor{G.node[B]}{state=done}
\recolor{G.node[C]}{state=done}
\recolor{G.node[D]}{state=done}
\recolor{G.node[T]}{state=done}
\recolor{G.edge[(S,A)]}{state=error}
\recolor{G.edge[(S,B)]}{state=error}
\recolor{G.edge[(A,B)]}{state=done}
\recolor{G.edge[(A,C)]}{state=done}
\recolor{G.edge[(B,D)]}{state=done}
\recolor{G.edge[(C,D)]}{state=done}
\recolor{G.edge[(C,T)]}{state=done}
\recolor{G.edge[(D,T)]}{state=done}
\apply{vars.var[path]}{value="done"}
\apply{vars.var[bottleneck]}{value="done"}
\annotate{G.node[S]}{label="min-cut: S side", color=error}

\end{animation}

\subsection{Key Insight: Residual Graph}

The \textbf{residual graph} includes:
\begin{itemize}
\item \textbf{Forward edges}: remaining capacity $= c(u,v) - f(u,v)$
\item \textbf{Reverse edges}: allow ``undoing'' flow $= f(u,v)$
\end{itemize}

Iteration 4 used a reverse edge to reroute: $A \to C$ flow decreased from 7 to 6, redirecting 1 unit through $C \to D \to T$.

\subsection{Complexity}

\begin{itemize}
\item \textbf{Edmonds-Karp} (BFS): $O(VE^2)$
\item \textbf{Dinic's}: $O(V^2 E)$ --- standard for competitive programming
\item \textbf{Push-Relabel}: $O(V^2 E)$ or $O(V^3)$
\end{itemize}

\subsection{Implementation (Dinic's)}

\begin{lstlisting}[language=Python]
from collections import deque

class Dinic:
    def __init__(self, n):
        self.n = n
        self.graph = [[] for _ in range(n)]

    def add_edge(self, u, v, cap):
        self.graph[u].append([v, cap, len(self.graph[v])])
        self.graph[v].append([u, 0, len(self.graph[u]) - 1])

    def bfs(self, s, t):
        self.level = [-1] * self.n
        self.level[s] = 0
        q = deque([s])
        while q:
            u = q.popleft()
            for v, cap, _ in self.graph[u]:
                if cap > 0 and self.level[v] < 0:
                    self.level[v] = self.level[u] + 1
                    q.append(v)
        return self.level[t] >= 0

    def dfs(self, u, t, f):
        if u == t:
            return f
        while self.iter[u] < len(self.graph[u]):
            e = self.graph[u][self.iter[u]]
            v, cap, rev = e
            if cap > 0 and self.level[v] == self.level[u] + 1:
                d = self.dfs(v, t, min(f, cap))
                if d > 0:
                    e[1] -= d
                    self.graph[v][rev][1] += d
                    return d
            self.iter[u] += 1
        return 0

    def max_flow(self, s, t):
        flow = 0
        while self.bfs(s, t):
            self.iter = [0] * self.n
            while True:
                f = self.dfs(s, t, float('inf'))
                if f == 0:
                    break
                flow += f
        return flow
\end{lstlisting}
